% !PRINTED 238
\weberpage{238}
%<<CONT
so erkennt man, daß $j(\omega)$ sich in eine nach Potenzen von $q^2$
fortschreitende Reihe entwickeln läßt, welche die Form hat
\weq{4}{j(\omega) = q^{-2} + a_1 + a_2 q^2 + a_3 q^4 + \cdots,}
worin die $a_1$, $a_2$, $a_3 \ldots$ \gesperrt{ganze rationale Zahlen} sind, die sich
successive berechnen lassen (es ergibt sich z. B. $a_1 = 744$,
$a_2 = 196\,884$). Die Entwickelungen der Größen (1) beginnen
also mit
\weq{5}{e^{-\frac{2\pi i c}{a}} e^{-\frac{2\pi i \partial}{a}\omega}}
\gesperrt{und sind daher alle voneinander verschieden. Die
Invariantengleichung ist also nach § 65 irreducibel.}

Es gibt aber einen zweiten Weg, um zu dieser wie überhaupt
zu den Transformationsgleichungen zu gelangen, den wir jetzt
zunächst bei der Invariantengleichung kennen lernen wollen.
Diese Ableitung stützt sich auf die Sätze des § 54 über Modulfunktionen
und gilt auch für ein \gesperrt{gerades} $n$. Hier ist es der
Satz 4, § 54, der zur Anwendung kommt:

I. \gesperrt{Jede Modulfunktion, die durch die beiden Transformationen}
\weq{6}{\begin{pmatrix} 1, & 0 \\ 1, & 1 \end{pmatrix},}
\weq{7}{\begin{pmatrix} 0, & 1 \\ -1, & 0 \end{pmatrix}}
\gesperrt{ungeändert bleibt, ist eine rationale Funktion von $j(\omega)$.}

Wir weisen zunächst nach, daß durch Anwendung der Substitutionen
(6), (7) die $\nu$ Größen (1) untereinander vertauscht
werden. Wir haben die Zusammensetzung
\weq{8}{\begin{pmatrix} a, & 0 \\ c, & \partial \end{pmatrix}\begin{pmatrix} 1, & 0 \\ 1, & 1 \end{pmatrix} = \begin{pmatrix} 1, & 0 \\ \lambda, & 1 \end{pmatrix}\begin{pmatrix} a, & 0 \\ c_1, & \partial \end{pmatrix},}
wenn
\weq{9}{c_1 = c + \partial - \lambda a}
gesetzt wird, und es geht durch die Transformation (6), da $j(\omega)$
durch jede lineare Transformation ungeändert bleibt,
\[
j\left(\frac{c+\partial\omega}{a}\right) \text{ in } j\left(\frac{c_1+\partial\omega}{a}\right)
\]
über.

Ferner bestimmen wir $a_2$, $\partial_2$, $c_2$, so daß
\weq{10}{\begin{gathered}
\begin{pmatrix} a, & 0 \\ c, & \partial \end{pmatrix}\begin{pmatrix} 0, & 1 \\ -1, & 0 \end{pmatrix} = \begin{pmatrix} \alpha, & \beta \\ \gamma, & \delta \end{pmatrix}\begin{pmatrix} a_2, & 0 \\ c_2, & \partial_2 \end{pmatrix}, \\
\alpha\delta - \beta\gamma = 1.
\end{gathered}}
